% BHU PYQ -- Professional Exam Template
\documentclass[11pt,a4paper]{article}

\usepackage[a4paper,top=2.5cm,bottom=2.5cm,left=2.8cm,right=2.8cm,headheight=14pt]{geometry}
\usepackage{amsmath,amssymb}
\usepackage[T1]{fontenc}
\usepackage[utf8]{inputenc}
\usepackage{lmodern}
\usepackage{microtype}
\usepackage{fancyhdr}
\usepackage{enumitem}
\usepackage{listings}
\usepackage{hyperref}
\usepackage{lastpage}
\usepackage{parskip}

\hypersetup{colorlinks=true,urlcolor=black,linkcolor=black,
  pdftitle={BVCA-302: Robotics Arduino -- BHU PYQ}}

\pagestyle{fancy}
\fancyhf{}
\renewcommand{\headrulewidth}{0.4pt}
\renewcommand{\footrulewidth}{0.4pt}
\fancyhead[L]{\small\textit{Banaras Hindu University}}
\fancyhead[R]{\small\textit{BVCA-302: Robotics Arduino}}
\fancyfoot[C]{\small Page \thepage\ of \pageref{LastPage}}

\lstset{
  basicstyle=\small\ttfamily,
  frame=single,
  breaklines=true,
  columns=flexible,
  keepspaces=true
}

\newlist{parts}{enumerate}{2}
\setlist[parts,1]{label=(\alph*),leftmargin=*,itemsep=4pt,topsep=3pt}
\setlist[parts,2]{label=(\roman*),leftmargin=*,itemsep=2pt,topsep=2pt}
\newcommand{\pts}[1]{\hfill\textit{[#1]}}

\begin{document}
\begin{center}
  {\large\bfseries BANARAS HINDU UNIVERSITY}\\[2pt]
  {\normalsize B.Voc. Semester III Examination, 2023-24}\\[6pt]
  \rule{0.8\linewidth}{0.5pt}\\[6pt]
  {\large\bfseries Paper: BVCA-302 --- Robotics Arduino}\\[4pt]
  {\normalsize Time: Three hours \hfill Maximum Marks: 70}
\end{center}

\rule{\linewidth}{0.4pt}

\smallskip
\noindent\textit{Instructions: | 1. Attempt any two out of four questions in Section A.}

\bigskip

- BGI INO. cia dhatjecssissosncoeditarcenrtpotterse
BVCA-302 : Robotics with Arduino
ee ee
| Instructions:
| 1. Attempt any two out of four questions in Section A. (2*12=24 Marks)

\medskip\noindent\textbf{Question 2.} Attempt any six out of nine questions in Section B (6*6=36 Marks)

\medskip\noindent\textbf{Question 3.} Attempt all 10 questions in Section C (1*10=10 Marks )

\bigskip\noindent{\large\bfseries SECTION --- A}
\rule{\linewidth}{0.4pt}\smallskip

(Long Answer Type questions to be answered in about 500 words each)

\medskip\noindent\textbf{Question 1.} Explain the fundamental components and functions of both hydraulic and
| pneumatic systems used in robotics. Provide a detailed comparison between the
two systems, highlighting their similarities and differences. Additionally, discuss
- the advantages and limitations of each system in the context of robotics
applications. :
| 2) Discuss the components and operation of a machine vision system used in
| robotics. Explain the role of lighting, cameras, frame grabbers, and computer
= ‘ processors in capturing and processing visual data for robot guidance and control.
| Provide examples of machine vision applications in industrial automation and
| robotics.
‘ 3) Explain the basic components of a hydraulic system and their functions. Provide
examples of applications where hydraulic systems are commonly used.
| 4) Describe the operation of a pneumatic system, including its basic components
and their functions. Discuss how a compressor generates compressed air, and
how pneumatic actuators convert this energy into mechanical power.
| : SECTION-B
| 5. Answer any six of the following in about 200 words each:
: (a) Define the term "robotics" and explain its significance in modern industrial
automation. Provide examples of industries where robotics is commonly utilized,
highlighting its impact on productivity and quality improvement.
\begin{parts}
  \item Explain the concept of degrees of freedom in industrial robots. Describe how
  degrees of freedom are related to the robot's ability to perform various tasks and
  navigate its work envelope efficiently.
  \item Explain the historical development of robotics from its origins in science fiction
  to its modern applications in industries. Highlight key milestones and
  technological advancements that have shaped the field.
  \item Discuss the advantages and disadvantages of different drive systems used in .
  industrial robots, including hydraulic drive, electric drive, and pneumatic drive.
  Provide examples of scenarios where each drive system is preferred based on
  performance requirements.
  \item Define actuators and sensors in the context of robotics. Explain their respective
  roles in robotic systems and provide examples of each. Discuss how actuators
  and sensors work together to enable robotic motion and perception in industrial
  and non-industrial applications.
  \item Discuss the importance of precision in robot movement and its components:
  P.T. 0.
  spatial resolution, accuracy, and repeatability.
  \item (c) Differentiate between fixed automation, programmable automation, and flexible
  automation. Provide examples of industries where each type of automation is
  commonly employed, and discuss their respective advantages and limitations. |
  \item Explain the concept of ROS (Robot Operating System) and its significance in
  robotics development.
  \item Describe the various types of motors used in robotics, including DC motors, AC
  motors, stepper motors, and servomotors.
\end{parts}

\bigskip\noindent{\large\bfseries SECTION --- C}
\rule{\linewidth}{0.4pt}\smallskip

\medskip\noindent\textbf{Question 6.} Answer the following objective type questions:
\begin{parts}
  \item In robot motion planning, what is the purpose of the "configuration space"?
  \item The physical space occupied by the robot
  \item The space of all possible joint configurations of the robot
  \item The workspace where the robot performs tasks
  \item The space-time continuum in which the robot operates
  \begin{parts}
    \item What is the role of the "end effector" in a robot manipulator?
    ‘ a) To control the robot's motion
  \end{parts}
  \item To provide feedback to the operator
  \item To act as the hand or tool of the robot
  \item To process sensor data
  \begin{parts}
    \item Which type of robot programming involves specifying the desired end-effector
    positions and orientations without detailing the intermediate joint motions?
  \end{parts}
  \item Motion Planning
  \_ b) Trajectory Planning
  \item Inverse Kinematics
  \item Forward Kinematics
  \begin{parts}
    \item What does the term "ROS package" refer to in the Robot Operating System
    (ROS)?
  \end{parts}
  \item A physical package containing robot components
  \item A collection of software libraries, tools, and configuration files
  \item The payload carried by a robot
  \item A unit of measurement for robot speed :
  (Vv) Which concept in robotics involves a robot's ability to recover from errors or
  unexpected events without human intervention?
  \item Fault Tolerance
  \item Redundancy
  \item Closed-Loop Control
  \item Safety Interlock
  \begin{parts}
    \item In the context of swarm robotics, what does the term "stigmergy" refer to?
  \end{parts}
  \item Communication through direct signalling
  \item Cooperation based on the use of pheromones or environmental cues
  \item Hierarchical control structure
  \item Swarm navigation using GPS
  | (vii) What is the origin of the term "robot"?
  . wine Count - - @
  >)
  \item Greek
  \item Czechoslovakia
  \item German
  \begin{parts}
    \item When did the modern technology of industrial robotics begin?
  \end{parts}
  \item 1920
  \item 1960
  \item 1980
  \item 2000
  \begin{parts}
    \item How is robotics defined as industrial automation?
  \end{parts}
  \item A subset of automation
  \item Unrelated to automation
  \item Only involves computer programming
  \begin{parts}
    \item What is another name for the Cartesian coordinate configurable robot?
  \end{parts}
  \item Spherical robot
   b) XYZ robot :
  \item Rectangular robot
  \item Hyperbolic robot
  = RRR EERE KKK KER EE
\end{parts}

\vfill
\rule{\linewidth}{0.4pt}
\begin{center}\small
\href{https://bhu-exam.vercel.app/}{Click here to visit our website}\\[4pt]\href{https://me-aryan.vercel.app/}{Click here to visit our contributor}\\[4pt]\href{https://quashan-me.vercel.app/}{Click here to visit our contributor}
\end{center}
\end{document}
